\section{一致可积性阻止期望质量逃逸}
\label{sec:12-Real-Analysis-Support/05-Measure-Integration-and-Conditional-Expectation/06}

一条离线评估流水线每天跑一次：用旧策略的日志，乘上第四节那种重要性权重，估计新策略的平均收益。多数日子给出的数字在小数点后第二位上下浮动，看着很稳。偶尔某一天，数字跳到平时的三倍，回查发现是几十条样本拿到了极大的权重。样本量加倍，这种日子照样出现。

权重的期望是 1，理论上估计量无偏，收敛也该有。可实际表现像是均值不存在。

问题不在收敛本身，在于收敛的方式与期望的算法之间有一道缝。几乎处处收敛逐条盯样本路径，依概率收敛盯偏差事件的概率，两者都不管坏事件发生时数值有多大。而期望把概率与幅度乘在一起。概率缩得再快，幅度涨得更快，乘积就不掉。

\subsection{一个消失却不掉质量的尖峰}

在 \(\Omega=(0,1)\) 上取均匀概率，令
\[
X_n=n\,\ind_{(0,1/n)}.
\]

\begin{center}
\begin{tikzpicture}[>=stealth,x=4cm,y=0.55cm]
\draw[gray!70] (0,0) rectangle (1,1);
\draw[gray!70] (0,0) rectangle (0.5,2);
\draw[thick] (0,0) rectangle (0.25,4);
\draw[->] (-0.03,0) -- (1.12,0);
\draw[->] (0,-0.3) -- (0,4.9);
\node[gray] at (0.75,0.55) {\footnotesize \(n=1\)};
\node[gray] at (0.375,1.5) {\footnotesize \(n=2\)};
\node at (0.4,3.9) {\footnotesize \(n=4\)};
\node at (0.63,4.3) {\footnotesize 每块面积都是 \(1\)};
\fill (0.62,0) circle (0.6pt);
\node at (0.62,-0.75) {\footnotesize 固定的 \(\omega\)};
\node at (0.62,-1.6) {\footnotesize 早晚落在尖峰之外};
\node at (1.06,-0.75) {\footnotesize \(\omega\)};
\end{tikzpicture}
\end{center}

\emph{尖峰越长越高、越缩越窄，每一条路径终将看不到它，可它带走的期望质量一点没少。}

对固定的 \(\omega>0\)，只要 \(n>1/\omega\) 就有 \(X_n(\omega)=0\)，所以 \(X_n\to 0\) 几乎处处，也依概率收敛。但
\[
\E[X_n]=n\cdot P\bigl((0,1/n)\bigr)=1
\]
对每个 \(n\) 都成立，\(\E[\abs{X_n-0}]=1\) 同样不动。极限的期望是 0，期望的极限是 1。

看清质量往哪儿跑了：它没有消散，而是挤进一个越来越窄、越来越高的地方。无论你把观察阈值 \(K\) 设到多大，总有靠后的 \(n\) 让 \(X_n\) 的全部质量藏在阈值之上。评估流水线里那些偶发的大数字，是同一件事的实测版本——权重的质量集中在极稀有的样本上，而那些样本平时不出现。

\subsection{一致可积性把阈值提到整族外面}

\begin{definition}
可积随机变量族 \(\mathcal X\) 称为\textbf{一致可积}的，若
\[
\lim_{K\to\infty}\ \sup_{X\in\mathcal X}\ \E\bigl[\abs{X}\,\ind_{\set{\abs{X}>K}}\bigr]=0.
\]
\end{definition}

量词的顺序是这条定义的全部要害：先取一个共同阈值 \(K\)，再对整族取最坏情况，最后让 \(K\) 变大。它要求存在一条对所有成员都管用的线，线以上的部分携带的期望质量可以一致地压到任意小。

顺序反过来就没有内容。每个单独的可积变量当然都有自己的阈值——可积的定义直接给出 \(\E[\abs X\ind_{\set{\abs X>K}}]\to 0\)——但这些阈值可以随 \(n\) 一起逃向无穷。尖峰族正是如此：给定任意 \(K\)，取 \(n>K\)，则 \(\set{X_n>K}=(0,1/n)\)，而
\[
\E\bigl[X_n\ind_{\set{X_n>K}}\bigr]=1.
\]
上确界卡在 1，怎么抬 \(K\) 都不降。这一族不一致可积。

一致可积推出 \(L^1\) 一致有界，反过来不成立：尖峰族的 \(\E[\abs{X_n}]\) 恒为 1，一致有界得很，却让质量跑掉了。总质量有界只说明账面上的钱没变多，不妨碍它全部集中到一个越来越难遇上的角落。

\(p>1\) 时的矩界就够用了。若 \(\sup_{X\in\mathcal X}\E[\abs X^{p}]\le C\)，则在 \(\set{\abs X>K}\) 上有 \(\abs X\le \abs X^{p}/K^{p-1}\)，于是
\[
\sup_{X\in\mathcal X}\E\bigl[\abs X\ind_{\set{\abs X>K}}\bigr]\le\frac{C}{K^{p-1}}\longrightarrow 0.
\]
多出来的那一点指数 \(p-1\)，就是压住尾部的全部本钱。这也解释了为什么重要性权重只有 \(\E_Q[W]=1\) 远远不够，而 \(\E_Q[W^{2}]<\infty\) 通常够——后者正是实践中``有效样本量''那个指标在量的东西。

控制收敛定理里那张共同包络 \(g\) 也推出一致可积，见 \hyperref[sec:12-Real-Analysis-Support/04-Limits-and-Integration/02]{``控制收敛用一个可积包络守住质量''}。反向不真：有些族一致可积，却找不到自然的逐点包络。一致可积因此是更宽的条件，包络是更容易验证的条件。

\subsection{Vitali 定理补上收敛图里的那个缺口}

有了这个条件，收敛之间的升级关系就补齐了。

\begin{theorem}[Vitali]
若 \(X_n\xrightarrow{\;P\;}X\) 且 \(\set{X_n}\) 一致可积，则 \(X\) 可积，且 \(\E[\abs{X_n-X}]\to 0\)，因而 \(\E[X_n]\to\E[X]\)。
\end{theorem}

骨架是把误差劈成两段：阈值以内与阈值以外。以内那段有统一上界，依概率收敛足以把它压小；以外那段由一致可积性统一压小，与 \(n\) 无关。两个条件各管一半——依概率收敛管住绝大多数样本的位置，一致可积管住少数极端样本的幅度。缺哪一半都不行：尖峰族满足前者不满足后者，结论就假。

控制收敛定理由此成为 Vitali 定理的一个易验证特例。各种收敛之间的关系见 \hyperref[sec:05-Probability/08-Limit-Theorems/01]{``随机变量的多种收敛''}；一致可积恰是从依概率收敛跨到 \(L^1\) 收敛所缺的那一步。

\subsection{停时把这个条件推到台前}

鞅在每个固定时刻都保持期望：\(\E[M_t]=\E[M_0]\)。人们想把它推到一个随机的停止时刻 \(\tau\) 上，得到 \(\E[M_\tau]=\E[M_0]\)。这个推广不是自动的。

赌徒倍投是标准反例：每输一次就把赌注翻倍，直到第一次赢了收手。这个停时几乎必然有限，停下来时必赚一个单位，于是 \(\E[M_\tau]=1\ne 0=\E[M_0]\)。破绽在停时之前的资产 \(M_{t\wedge\tau}\) 上：连输 \(k\) 次的概率是 \(2^{-k}\)，此时负债是 \(2^{k}-1\)，概率与幅度恰好抵消。这与本节开头那张图是同一个形状——越来越罕见、越来越极端，质量逃到那里去了。

可选停时定理成立需要的补丁正是一致可积：若族 \(\set{M_{t\wedge\tau}}_{t\ge 0}\) 一致可积（有界停时、有界鞅、或者有共同的可积包络都能保证），则 \(\E[M_\tau]=\E[M_0]\)。细节见 \hyperref[sec:06-Random-Processes/04-Applications/02]{``Martingale：公平游戏与停时''}。上一节末尾留的那个口子——\(M_t=\E[Y\given\mathcal F_t]\) 是否收敛到 \(Y\)——答案也在这里：这类鞅自动一致可积，收敛得以在 \(L^1\) 意义下成立。

\subsection{它在估计与训练里的常见形态}

重要性权重是最直接的一处。支撑覆盖回答``能不能估''，一致可积回答``估出来的东西稳不稳、极限与期望能不能交换''。前者靠改采样分布解决，后者靠截断权重、自归一化、或者干脆换一个方差更小的估计量解决——截断权重引入偏差换回一致可积，这笔交易在实践中通常划算。

同一形态在别处反复出现。带重尾的收益分布下样本均值收敛得极慢，见 \hyperref[sec:05-Probability/06-Common-Distributions/07]{``重尾分布：当均值与方差失效时''}；Monte Carlo 估计的方差诊断，见 \hyperref[sec:05-Probability/08-Limit-Theorems/05]{``Monte Carlo：从随机样本到数值答案''}；训练中的梯度估计遇上偶发的巨大梯度，梯度裁剪做的事与截断权重同构。它们共享一句诊断：均值看着存在，估计却不稳，多半是质量正在往罕见而极端的区域集中。

所以看到 \(X_n\to X\) 之后，别急着在两边加期望。先问收敛是哪一种，再找共同包络、更高阶矩界或一致可积性其中之一。极限存在与平均质量随极限保存，是两个结论，中间隔着这一整节。

这也是整个单元的收束。第一节划定哪些问题可问，第二节给它们定大小，第三节把大小加起来，第四节说清每个密度相对于谁，第五节把``利用信息''写成投影，而这一节交代了：当极限出现时，这套账还需要一个条件才能继续对得上。
